\section{Eastern Christian transmission: several books, not one rite}

The history becomes geographically broader when Greek, Armenian, West-Syriac,
and Coptic witnesses are placed beside the Latin books. It also becomes less
linear. These witnesses do not disclose an undifferentiated ``Eastern
exorcism.'' They preserve different kinds of textual objects: prayers gathered
in a Byzantine euchologion, catechumenal material embedded in baptismal orders,
and later editions that compare several manuscript families. Their value lies
first in the architecture of the books and in the distinctions that architecture
requires.

This section is historical, not juridical or operational. It does not state the
present law of an Eastern Catholic Church or an Orthodox Church, identify a
currently authorized minister, reproduce a formula, or offer a sequence for
use. The words \emph{Byzantine}, \emph{Armenian}, \emph{West-Syriac}, and
\emph{Coptic} identify bounded textual and manuscript settings here. They do
not erase differences of communion, jurisdiction, recension, language,
territory, or date. In particular, a historical Armenian or Syriac witness
cannot be assigned retrospectively to one present Catholic or Orthodox body
without evidence of its provenance and reception.

\subsection{A Byzantine prayer-book stream}

The Byzantine evidence includes prayers transmitted as units in euchologia.
That location differs from a clergy list, an ordination form, a baptismal
scrutiny, and the later Latin section for persons described as afflicted. The
euchologion is a service book containing many kinds of ecclesial prayer. The
coexistence of material inside it establishes a book environment; it does not
by itself establish how commonly each item was used, by whom, or under what
local authorization.

The Barberini Greek manuscript now numbered Vatican, Barberini gr.~336 supplies
an especially important control. André Jacob's codicological description maps
a group at folios 84r--95r. Within that span he identifies a Basil-attributed
unit at folios 84r--88v, a Chrysostom-attributed unit at folios 88v--91r,
additional prayers at folios 91r--95r, and a lacuna within the sequence. The
physical map matters more than the famous names. Two prayers later printed
under an attribution to John Chrysostom occur anonymously in this early
witness. Modern editors Stefano Parenti and Elena Velkovska print the bounded
Barberini material at pages 225--227 of their 1995 edition; Miguel Arranz gives
a further Constantinopolitan euchologion control at pages 371--372 of his 1996
edition. These protected modern editions are bibliographic controls only here;
their text is neither retained nor reproduced.

This evidence imposes a strict authorship ceiling. A prayer attributed in later
transmission to Basil or Chrysostom is not thereby proved to have been composed
by that Father. The manuscript heading, the presence or absence of a name, the
date of the physical witness, and the date of composition are separate facts.
Even the oldest controlled witness supplies a terminus before which the item
was copied, not a biography of its author or a complete history of its earlier
forms.

Jacques Goar's \emph{Euchologion sive rituale Graecorum} makes the later printed
collection directly visible. The public-domain seventeenth-century edition and
the expanded 1730 printing place the relevant units together; the latter's
printed pages 578--584 became a common reference point for later cataloguers
and editors. Goar is valuable as evidence for early-modern editorial
consolidation and for locating the units that later scholarship compares. He
is not a substitute for the manuscript. His arrangement cannot be projected
backward unchanged into the eighth century, and his headings cannot settle the
authorship of anonymous or variably attributed material.

Dionysios Stathakopoulos's manuscript-history control sharpens this point. He
distinguishes the Barberini witness from an eleventh-century Zürich booklet,
an Athenian manuscript dated 1710, and an Athonite manuscript dated 1735.
Those books transmit different combinations of attributed and anonymous
material. The Zürich witness includes units attributed to several Fathers; the
later Athenian and Athonite books add other names and forms. Their plurality is
not noise to be harmonized away. It is positive evidence that the corpus was
open to selection, attribution, rearrangement, and local transmission.

The safe Byzantine claim is therefore modest but consequential. Identified
Greek euchologia preserve bounded prayer units concerning persons described
as troubled by unclean spirits, and some of those units entered an influential
early-modern printed collection. The manuscript tradition is older and more
varied than Goar's arrangement. It does not establish patristic authorship,
uniform performance, ordinary frequency, efficacy, a single Byzantine
procedure, or a present Orthodox rule.

The survival pattern also warns against treating the most accessible edition
as the most representative witness. Goar's large printed volume is easier to
consult than a manuscript folio and later became a standard bibliographic
coordinate. That practical prominence can conceal the asymmetry of the
evidence. One manuscript may preserve only part of a later group; a heading may
be absent; a codex may have lost leaves; and a later copy may join material
that circulated separately. Historical synthesis must retain those
asymmetries. A table that lists seven later-attributed prayers, for example,
would describe one collected state. It would not prove that every earlier
euchologion contained seven, that all seven had the same history, or that one
community used them as a fixed set.

Nor does the label \emph{euchologion} solve the question of use. It identifies
a class of book, not an attendance register or pastoral report. A copied item
may have been intended for use, preserved for reference, reproduced from a
model, or transmitted because the compiler valued completeness. Without
independent evidence, the manuscript cannot decide among those possibilities.
This is why the present account speaks of transmission and placement rather
than incidence or outcome.

\subsection{Armenian initiation in identified manuscripts}

The Armenian comparison changes the type of evidence. Frederick Cornwallis
Conybeare's 1905 \emph{Rituale Armenorum} prints and translates a baptismal
order from identified Armenian manuscripts, chiefly Barberini Oriental
vii.44 and Vienna Mechitarist codex 68, while recording variants and omissions
in other witnesses. The relevant unit occupies printed pages 86--93. It moves
from rites associated with the catechumenate into the canon of baptism, with
prayers over the candidate, a renunciation-and-profession boundary, and entry
into the baptismal sequence.

The editorial notes are essential evidence. Conybeare repeatedly indicates
that one manuscript excludes a rite, another changes a rubric, and other
witnesses expand or abbreviate the sequence. The printed text is therefore not
a transparent photograph of ``the Armenian rite.'' It is an old scholarly
edition of a manuscript family whose agreements and differences are exposed on
the page. Its English translation provides lawful access to the architecture,
but it remains Conybeare's translation and editorial construction.

This witness also protects a distinction already visible in Jerusalem and
Western initiatory sources. Exorcistic and renunciatory language may belong to
preparation for Christian initiation. Its recipient and purpose are not
therefore identical with the recipient and purpose of a later rite directed
toward a person alleged to be possessed. In Conybeare's bounded pages, the
material is ordered toward baptismal incorporation. To extract one expression
from that setting and label it an Armenian major exorcism would destroy the
book's own category.

The historical coordinate ``Armenian'' requires another restraint. The
manuscripts witness an Armenian liturgical tradition, but this section has not
established their use as present law in the Armenian Apostolic Church, the
Armenian Catholic Church, or any eparchy. Nor does manuscript agreement prove
actual frequency. The safe claim is that identified Armenian manuscripts,
known through a public-domain edition, transmit differentiated catechumenal and
baptismal preparation. They broaden the history of Christian exorcistic
language while refusing the fiction of one Eastern procedure.

\subsection{West-Syriac orders and the problem of names}

Heinrich Denzinger's \emph{Ritus Orientalium}, volume I, offers a further
nineteenth-century window. It gathers Latin translations and editorial notes
for Coptic, Syriac, and Armenian sacramental orders from earlier collections
and manuscripts. Its age and derivative character lower its authority for
exact original-language wording, but its explicit source notes and parallel
orders make it useful for bounded historical comparison.

One West-Syriac baptismal order begins at printed page 266 and places its
exorcistic and renunciatory preparation at pages 272--273. Another order,
printed under the name of Jacob of Sarug, begins at page 334 and places
comparable preparatory material from page 337 onward. Elsewhere Denzinger
discusses orders associated in the received tradition with Jacob of Edessa,
Severus of Antioch, and Basil. His own prolegomena expose uncertainty about
relationships among those orders and about the labels assigned by compilers.

The names must consequently be treated as transmission labels, not as
authenticated signatures. ``Of Jacob,'' ``of Severus,'' or ``of Basil'' may
record an attribution, an association, an editorial title, or a claim received
from a source used by Denzinger. It does not establish composition by that
person. The same restraint applied to the Greek Fathers applies here with
greater force because the consulted edition mediates Syriac material through
Latin and through earlier editors.

What the pages do establish is structural. More than one identified
West-Syriac baptismal order places exorcistic material within preparation for
initiation. Their internal differences prevent a single composite text from
being reconstructed by taking whatever is fullest from each. They also prevent
an argument from silence: omission in one abbreviated order does not prove
absence from a whole community, while presence in one manuscript does not
prove universal use.

The ecclesial boundary remains equally important. The historical descriptor
``West-Syriac'' does not decide whether a specific witness belongs to the
ancestry or use of the Syriac Orthodox Church, a Syriac Catholic Church, the
Maronite Church, or another community at every point in its transmission.
Denzinger sometimes writes with nineteenth-century confessional terminology
that must itself be historicized. A present-law claim would require a named
Church, its competent authority, an approved book, an exact edition and
language, a territory, and a date. None is supplied by this old comparative
edition.

\subsection{A Coptic initiatory comparison}

Denzinger's second Coptic baptismal order supplies a fourth regional control at
printed pages 215--219. The unit places prayers over catechumens and preparatory
anointing before a renunciation-and-profession boundary at page 217, then
continues into baptism. As with the Armenian and West-Syriac witnesses, the
governing architecture is initiatory. The material belongs to the making of a
catechumen and the approach to the font, not to a free-standing manual for
testing an afflicted person.

This control is deliberately narrower than the word ``Coptic'' might suggest.
Denzinger prints Latin translations and cites earlier learned collections. His
edition does not give this study a newly checked Coptic text, a manuscript
apparatus, or a secure recension history. The separate 1848 Tattam collection
previously probed by this project combines Coptic and Arabic church-order
material whose constituent boundaries and relationships have not been
resolved. It remains a lead rather than independent support for a generic
claim about Coptic practice.

The Denzinger pages can nevertheless support one useful conclusion. Coptic
initiatory evidence cannot simply be replaced by the Greek euchologion or by a
Latin ritual. It has its own editorially transmitted order and must be studied
through its own languages, manuscripts, and ecclesial history. The current
Coptic Orthodox and Coptic Catholic Churches are not interchangeable
authorities, and neither one's present discipline is established by an 1863
Latin compilation.

\subsection{Three editorial layers}

The controls themselves belong to three different scholarly moments.
Seventeenth-century Goar presents a Greek--Latin service-book collection.
Nineteenth- and early-twentieth-century Denzinger and Conybeare organize and
translate eastern sacramental material for learned readers. Jacob,
Stathakopoulos, and the modern critical editors return attention to the
codices, their foliation, lacunae, anonymous headings, and changing
combinations. These layers cannot be treated as interchangeable quotations
from an abstract tradition.

For this study the public-domain editions supply reproducible page-level
controls. Their exact scans or text layers were identified and hashed, and the
reader claims stay at the level that those editions can carry. Modern
scholarship supplies the correction that old editorial titles and
attributions require manuscript-level testing. Because the modern works remain
copyrighted, this project records their bibliographic identity and bounded
loci without retaining or reproducing their protected text.

The layered method leaves genuine work open. A full account would collate
original Greek, Armenian, Syriac, Coptic, and Arabic witnesses; distinguish
recensions; test dates and provenance; and investigate reception in named
communities. It would also separate historical book use from current approved
books and particular law. The absence of that larger project is not repaired
by silently promoting an old Latin translation into an original-language
critical edition. Instead, each present claim stops at the exact artifact that
was inspected.

\subsection{What comparison can and cannot show}

Placed together, the four streams produce a map of differences:

\begin{threestable}{Witness family}{What the checked control shows}{Claim
ceiling}
\textbf{Byzantine Greek} & Stand-alone or grouped prayer units in identified
euchologia; variable anonymity and attribution; later consolidation in Goar. &
No secure patristic authorship, uniform rite, frequency, efficacy, or present
Orthodox law.\\
\textbf{Armenian} & Catechumenal and baptismal architecture in identified
manuscripts printed with variants by Conybeare. & No composite Armenian
procedure or present Armenian Apostolic or Catholic rule.\\
\textbf{West-Syriac} & Several baptismal orders in Denzinger's comparative
edition, with inherited personal attributions and visible variation. & No
authenticated eponymous authorship, original-language critical text, universal
practice, or current Syriac or Maronite rule.\\
\textbf{Coptic} & Catechumenal preparation and renunciation within one bounded
baptismal order in Denzinger. & No complete Coptic corpus, recension history,
generic practice, or present Coptic Orthodox or Catholic rule.\\
\end{threestable}

The comparison does not yield a common denominator that can be performed. It
yields historical categories. First, exorcistic vocabulary circulated both in
initiatory settings and in separately gathered prayers. Second, a famous name
in a heading is evidence of attribution, not automatically of authorship.
Third, a printed service book is a stage in transmission rather than the origin
of every item it contains. Fourth, textual survival is not a statistic of use.
Fifth, resemblance across languages does not by itself prove borrowing, common
descent, or equal juridical status.

This broader map corrects two opposite distortions. A Latin-only history might
mistake the development of the Roman office and ritual book for the whole
Christian history of exorcistic language. A flattened East--West comparison
might then repair that omission by inventing one timeless ``Eastern rite.''
The witnesses support neither account. They show several communities receiving,
editing, and locating related language in materially different books.

The result also clarifies the relationship between historical and present
claims. An Eastern Catholic Church shares Catholic faith and communion while
retaining its own legal and liturgical identity. An Orthodox Church is not
governed by the Latin Code or by a Catholic dicastery. Historical texts shared
or revered across these boundaries do not erase them. CCEO canons 1--2 prevent
the Latin canon from being treated as Eastern Catholic law; they do not fill
the resulting research gap with a universal Eastern rule. Historical
euchologia likewise cannot legislate for a present Orthodox jurisdiction.

For the purposes of this study, the positive conclusion is therefore one of
plural transmission. Greek prayer books, Armenian manuscript orders,
West-Syriac baptismal families, and a Coptic baptismal order all enlarge the
documented geography. They also strengthen the non-operational method: name the
book, language, witness, date, recipient, and purpose before comparing a text.
Where those coordinates are missing, restraint is not a failure to synthesize.
It is the historical conclusion demanded by the evidence.
That conclusion remains revisable when a better critical edition, a newly
catalogued manuscript, or an identified Church's authoritative present book
supplies evidence that the currently controlled artifacts cannot provide.
